\documentclass[10pt,a4paper]{amsart}
\usepackage[margin=19mm]{geometry}
\usepackage{amsmath,amssymb,mathtools}
\usepackage[scr=rsfs,cal=euler]{mathalfa}
\usepackage{xfrac}
\usepackage[unicode,hidelinks,bookmarks=false]{hyperref}
\newcommand{\ie}{i.e.\ }
\newcommand{\cf}{cf.\ }
\newcommand{\resp}{respectively\ }
\newcommand{\Subsection}{\subsection}
\providecommand{\nobreakdash}{\nobreak\hbox{-}}
\setlength{\emergencystretch}{2em}
\multlinegap0pt
\usepackage{bm}
\usepackage{bbm}
\usepackage{dsfont}
\usepackage{xspace}
\usepackage{cancel}
\usepackage{mathtools}
\usepackage{amsthm}
\makeatletter
\@ifundefined{lem}{\newtheorem{lem}{Lemma}}{}
\makeatother
\title[Numerical Approximation and Bifurcation Results for an Ellip]{Numerical Approximation and Bifurcation Results for an Elliptic Problem with Superlinear Subcritical Nonlinearity on the Boundary}
\author{Shalmali Bandyopadhyay, Thomas Lewis, Dustin Nichols}
\date{Source: arXiv:2509.08990v1 (CC BY 4.0)}
\begin{document}
\maketitle

\noindent\emph{Source and licence.} Numerical Approximation and Bifurcation Results for an Elliptic Problem with Superlinear Subcritical Nonlinearity on the Boundary. Version arXiv:2509.08990v1 (CC BY 4.0), distributed under the Creative Commons Attribution 4.0 International licence, as recorded in the arXiv licence metadata for this exact version and shown on the abstract page https://arxiv.org/abs/2509.08990v1. Full rights notice: licenses/arxiv-2509.08990-CC-BY-4.0.txt.\par\smallskip

\noindent\textit{Editorial scope.} Counted unit: the Lem whose source label is \texttt{thm:admissibile} in the article (source0.tex lines 238--246, source label \texttt{thm:admissibile}), delivered together with the article's own complete original proof of it (source0.tex lines 247--386, one \texttt{proof} environment of 140 lines). No mathematics is written, repaired or re-derived: statement and proof are the article's own text line for line. Required context retained uncounted: none. Every definition, display and label of those lines is retained unchanged. Omitted from this package and not redistributed: 1--237, 387--1228. Nothing inside a retained excerpt is omitted or altered: every retained body is delivered exactly as the article's lines are written, so no omission receipt of any kind accompanies this package. Citations: the retained excerpts carry no citation command at all, so no bibliography entry of the article is transcribed and no reference is redistributed. Reference adaptation: none was needed and none was made; every reference inside the delivered unit resolves inside it, so no unresolved reference or citation is produced. Printed slips kept literal: no printed word, symbol, hypothesis or number of the counted statement or of its proof was altered. Layout and class: the article's own class is \texttt{amsart} with packages including amsmath, amssymb, amsthm, amstext, enumerate, graphicx, amsfonts, hyperref, epsfig, epstopdf, caption, inputenc, fontenc, tcolorbox; the wrapper is the lane's pinned \texttt{amsart} editorial wrapper at 10pt on A4, which restates the theorem-like environments the retained text needs and adds the article's own macro definitions that the retained text needs (none), with extra wrapper packages (bm, bbm, dsfont, xspace, cancel, mathtools). The delivered numbering is the unit's own. Assets: no retained span contains a figure environment, a graphics-inclusion command, a PDF-page inclusion, a table or a TikZ picture; no image file is copied into this package, the package declares no asset dependency and no image file is redistributed with it. Licence: version arXiv:2509.08990v1 of this article is distributed under the Creative Commons Attribution 4.0 International licence, as recorded in the arXiv licence metadata for this exact version and shown on the abstract page https://arxiv.org/abs/2509.08990v1. A full rights notice is delivered at licenses/arxiv-2509.08990-CC-BY-4.0.txt. Characters: every retained excerpt is pure ASCII, so the delivered engine, XeLaTeX with its default Latin Modern text font, reports no missing character.
\begin{lem}\label{thm:admissibile}
There exists a nonnegative grid function $\widetilde{u}_h : \mathcal{T}_h \to \mathbb{R}$ such that 
\begin{equation}\label{d_sol_single_tilde}
    \left\lbrace\begin{array}{rcll}
    -\Delta_{h} \widetilde{u}_{h}+ \widetilde{u}_{h}&=&0 & \text{ in } \mathcal{T}_h \cap \Omega \,, \\
    \nabla_h^* \widetilde{u}_{h}  \cdot \widehat{n} - \lambda \widetilde{f}(\widetilde{u}_{h}  ) &=& 0 & \text{ on } \mathcal{T}_h \cap \widetilde{\partial \Omega} . 
    \end{array}\right.
\end{equation}    
\end{lem}

\begin{proof}
Let $M_K > 0$ be defined by $M_K := 2 N \frac{K}{h_*^2}$.  
Define the grid function $\overline{u}_h : \mathcal{T}_h \to \mathbb{R}$ by 
\begin{equation}
    \overline{u}_h(x_\alpha) := 
    \begin{cases}
        M_K & \text{if } x_\alpha \in \mathcal{T}_h \cap \Omega , \\ 
        M_K + K & \text{if } x_\alpha \in \mathcal{T}_h \cap \partial \Omega . 
    \end{cases}
\end{equation}
Define the subgrid $\mathring{\mathcal{T}}_h$ by 
\[
    \mathring{\mathcal{T}}_h := \left\{ x_\alpha \in \mathcal{T}_h \mid 
    x_{\alpha \pm \mathbf{e}_i} \notin \partial \Omega \text{ for all } i=1,2,\ldots,N \right\} . 
\]
Then, there holds 
\begin{align*}
    -\Delta_h \overline{u}_h(x_\alpha) + \overline{u}_h(x_\alpha) 
    & = 0 + M_K \\ 
    & = M_K \geq 0 
\end{align*}
for all $x_\alpha \in \mathring{\mathcal{T}}_h$.  
Choose $x_\alpha \in (\mathcal{T}_h \setminus \mathring{\mathcal{T}}_h) \cap \Omega$, 
and define 
\[
    E_i^\pm (x_\alpha) := 
    \begin{cases}
        1 & \text{if } x_{\alpha \pm {e}_i} \in \partial \Omega , \\ 
        0 & \text{otherwise} 
    \end{cases}
\]
for all $i = 1,2,\ldots,N$.  
Then 
\begin{align*}
    -\Delta_h \overline{u}_h(x_\alpha) + \overline{u}_h(x_\alpha)
    & = \sum_{i=1}^N E_i^+(x_\alpha) \frac{M_K - (M_K+K)}{h_i^2} + \sum_{i=1}^N E_i^-(x_\alpha) \frac{M_K - (M_K+K)}{h_i^2} 
        + M_K \\ 
    & \geq - 2N \frac{K}{h_*^2} + M_K = 0 . 
\end{align*}
Lastly, for $x_\alpha \in \mathcal{T}_h \cap \widetilde{\partial \Omega}$, there holds 
\begin{align*}
    \nabla_h^* \overline{u}_{h}  \cdot \widehat{n} - \lambda \widetilde{f}(\overline{u}_{h})
    & = \frac{K}{h_i} - \lambda \widetilde{f}(M_K+K) \\ 
    & \geq \frac{K}{h^*} - \lambda \frac{K}{\lambda h^*} = 0  
\end{align*}
for some $i \in \{1,2,\ldots,N\}$.  
Hence, $\overline{u}_h$ is a positive supersolution of \eqref{d_sol_single_tilde}.  
Next, we define the grid function $\underline{u}_h : \mathcal{T}_h \to \mathbb{R}$ by 
\begin{equation}
    \underline{u}_h(x_\alpha) := 0 
\end{equation}
for all $x_\alpha \in \mathcal{T}_h$.  
Then, 
\begin{align*}
    -\Delta_h \underline{u}_h(x_\alpha) + \underline{u}_h(x_\alpha) 
    & = 0 
\end{align*}
for all $x_\alpha \in \mathcal{T}_h \cap \Omega$ 
and 
\begin{align*}
    \nabla_h^* \underline{u}_{h}  \cdot \widehat{n} - \lambda \widetilde{f}(\underline{u}_{h})
    & = - \lambda \widetilde{f}(0) \\ 
    & \leq - \rho < 0 
\end{align*}
for all $x_\alpha \in \mathcal{T}_h \cap \partial \Omega$.  
Hence, $\underline{u}_h$ is a nonnegative subsolution of \eqref{d_sol_single_tilde} 
with $\underline{u}_h$ not an actual solution due to the cutoff $\rho$ for $\widetilde{f}$.  
Furthermore, $\underline{u}_h \leq \overline{u}_h$.  

Let $S(\mathcal{T}_h)$ denote the space of all grid functions, 
and define the nonempty convex subspace $\overline{\underline{S}}(\mathcal{T}_h)$ by 
\[
    \overline{\underline{S}}(\mathcal{T}_h) := \left\{ v_h : \mathcal{T}_h \to \mathbb{R} \mid 
    \underline{u}_h \leq v_h \leq \overline{u}_h \right\} . 
\]
We use the Schauder fixed point theorem to show that \eqref{d_sol_single_tilde} has a solution in the space 
$\overline{\underline{S}}(\mathcal{T}_h)$ from which the result follows.  
Thus, \eqref{d_sol_single_tilde} has a nonnegative solution.  
Define the mapping $\mathcal{M} : S(\mathcal{T}_h) \to S(\mathcal{T}_h)$ by 
$w_h := \mathcal{M} v_h$ if 
\begin{equation}\label{cM_mapping}
    \left\lbrace\begin{array}{rcll}
    -\Delta_{h} w_{h}+ w_{h}&=&0 & \mbox{ in } \mathcal{T}_h \cap \Omega \,, \\
    \nabla_h^* w_{h}  \cdot \widehat{n} &=& \lambda \widetilde{f}(v_{h}  ) & \mbox{ on } \mathcal{T}_h \cap \widetilde{\partial \Omega} 
    \end{array}\right.
\end{equation}
for $v_h \in S(\mathcal{T}_h)$. 
Clearly a fixed point of $\mathcal{M}$ is a solution to \eqref{d_sol_single_tilde}.  

Let $n = | \mathcal{T}_h \cap (\Omega \cup \widetilde{\partial \Omega}) |$. 
Then, there exists a matrix $A \in \mathbb{R}^{n \times n}$ and a vector $\vec{b} \in \mathbb{R}^n$ such that 
\eqref{cM_mapping} is equivalent to solving $A \vec{w} = \vec{b}$, 
where $\vec{w}$ is the vectorization of $w_h$ 
and $\vec{b}$ depends on $\vec{v}$, the vectorization of $v_h$.  
By the choice of the difference operator $\nabla_h^*$ used on the boundary 
and the form of the equations corresponding to $-\Delta_h w_h + w_h$, 
we have the matrix $A$ is a symmetric Z-matrix.  
By Gershgorin's circle theorem, all of the eigenvalues of $A$ are nonnegative.  
Since constant functions are not in the nullspace of $A$, we actually have the matrix must be 
nonsingular using properties of the corresponding discretization of Poisson's equation with Neumann boundary conditions.  
Therefore, $A$ is a nonsingular M-matrix.  
It follows that $w_h$ is uniquely defined by $v_h$ in \eqref{cM_mapping} and that the mapping $\mathcal{M}$ is well-defined.  

Suppose $v_h \in \overline{\underline{S}}(\mathcal{T}_h)$, and let $w_h = \mathcal{M} v_h$.  
We show $w_h \in \overline{\underline{S}}(\mathcal{T}_h)$ 
from which it follows that $\mathcal{M}$ has a fixed point in the space $\overline{\underline{S}}(\mathcal{T}_h)$.  
Observe that 
\begin{align*}
    -\Delta_h \underline{u}_h(x_\alpha) + \underline{u}_h(x_\alpha) 
    & = 0 
    = -\Delta_{h} w_{h}+ w_{h} 
    \leq -\Delta_h \overline{u}_h(x_\alpha) + \overline{u}_h(x_\alpha) 
\end{align*}
for all $x_\alpha \in \mathcal{T}_h \cap \Omega$ 
and 
\begin{align*}
    \nabla_h^* \underline{u}_{h}  \cdot \widehat{n} 
    & = 0 
    \leq \lambda \widetilde{f}(v_h) 
    = \nabla_h^* w_{h}  \cdot \widehat{n}  
\end{align*}
with
\begin{align*}
    \lambda \widetilde{f}(v_h) 
    & \leq \lambda \frac{K}{\lambda h^*} 
    \leq \frac{K}{h_i} 
    = \nabla_h^* \overline{u}_{h}  \cdot \widehat{n} 
\end{align*}
for some $i \in \{1,2,\ldots,N\}$ 
for all $x_\alpha \in \mathcal{T}_h \cap \widetilde{\partial \Omega}$.  
Letting $\underline{U}, \vec{w}, \overline{U} \in \mathbb{R}^n$ denote the vectorizations of 
$\underline{u}_h, w_h, \overline{u}_h$, respectively, 
it follows that 
\[
    A \underline{U} \leq A \vec{w} \leq A \overline{U}  
\]
using the natural partial ordering for vectors.  
Since $A$ is an M-matrix, there holds $\underline{U} \leq \vec{w} \leq \overline{U}$, 
and we can conclude $w_h \in \overline{\underline{S}}(\mathcal{T}_h)$.  
\end{proof}

\end{document}
